\documentclass[11pt,a4paper]{article}
\usepackage[T1]{fontenc}
\usepackage[utf8]{inputenc}
\usepackage{amsmath,amssymb,amsthm}
\usepackage[margin=1in]{geometry}
\usepackage[colorlinks=true,linkcolor=blue,urlcolor=blue]{hyperref}
\theoremstyle{plain}
\newtheorem{theorem}{Theorem}
\newtheorem{lemma}[theorem]{Lemma}
\newtheorem{proposition}[theorem]{Proposition}
\newtheorem{corollary}[theorem]{Corollary}
\theoremstyle{definition}
\newtheorem{remark}[theorem]{Remark}

\title{Recovering the unwritten recurrences of the table OEIS A237158}
\author{Adrian Perez Fontelles\\ \small Independent researcher}
\date{2 September 2026}
\begin{document}
\maketitle

\begin{abstract}
OEIS A237158 is a two-dimensional table $T(n,k)$ whose empirical recurrences are, for several of
its lines, given only as an ORDER --- ``k=2: [order 16]'' and the like --- with the
coefficients in a linked file rather than in the entry. Those conjectures can still be settled,
and the recurrences recovered. A line of the table is a fixed-width or fixed-height array
count, hence a walk count on finitely many vertices, hence satisfies some monic recurrence of
order at most that number; Berlekamp--Massey on twice as many exact terms returns the MINIMAL
one. Where its order equals the order the entry states --- and where the same holds after the
first few terms are dropped --- any recurrence of that order the line satisfies has a
characteristic polynomial that is a multiple of the minimal one and of equal degree, hence is
that polynomial. This note settles column 2, column 3.
\end{abstract}

\noindent\small 2020 Mathematics Subject Classification. 05A15, 11B37, 15A18.\normalsize

\section{The table and the unwritten conjectures}

OEIS A237158 is ``T(n,k) = Number of (n+1) X (k+1) 0..3 arrays with the maximum plus the minimum of every 2 X 2 subblock equal.''. It is read by antidiagonals and begins
\[
256, 1868, 1868, 14096, 37652, 14096, 111900, 762508,\ \dots
\]
The entry carries, among its empirical claims:
\begin{quote}\small
k=2: [order 16]\\
k=3: [order 38]
\end{quote}
As of the ``Last modified'' line on the live entry (Apr 14 2015 08:32:52, revision 8) these are still
recorded as empirical, and the recurrences themselves are not in the entry text.

\section{Why a line of the table is a walk count}

Fix the index that the line fixes. The entry defines $T(n,k)$ as a number of arrays of a shape
determined by $n$ and $k$, subject to a condition local to a bounded window of consecutive
lines. Substituting the fixed index into the entry's own wording turns the definition into an
ordinary one-parameter array count, and for such a count the admissible windows are the
vertices of a finite digraph and an array is a walk. So the line is a walk count on $S$
vertices, and by the Cayley--Hamilton theorem it satisfies a monic integer recurrence of order
at most $S$.

\section{Recovering the recurrences}

\begin{lemma}\label{lem:bm}
A sequence known to satisfy some linear recurrence of order at most $S$ is determined, with its
minimal recurrence, by its first $2S$ terms; Berlekamp--Massey applied to them returns that
minimal recurrence exactly.
\end{lemma}

\subsection*{Column $k=2$}
Substituting the fixed index into the entry's wording gives an ordinary one-parameter array
count, a walk on $S=136$ vertices. Berlekamp--Massey on $2S=272$ exact terms
returns a minimal recurrence of order $16$ --- the order the entry states --- with
integer coefficients
\begin{center}\small\begin{tabular}{cccccccc}
$c_{1}$ & $55$ & $c_{5}$ & $-418194$ & $c_{9}$ & $35087663$ & $c_{13}$ & $36204908$ \\
$c_{2}$ & $-1035$ & $c_{6}$ & $1564374$ & $c_{10}$ & $12182197$ & $c_{14}$ & $-9872064$ \\
$c_{3}$ & $6857$ & $c_{7}$ & $1495550$ & $c_{11}$ & $-64353607$ & $c_{15}$ & $-5608640$ \\
$c_{4}$ & $16166$ & $c_{8}$ & $-20565065$ & $c_{12}$ & $12508772$ & $c_{16}$ & $1752064$ \\
\end{tabular}\end{center}
so that $a(m)=\sum_{i=1}^{16}c_i\,a(m-i)$ for every $m$. Removing the first
$16+4$ terms and repeating gives order $16$ again, which is what the
identification below needs.

\subsection*{Column $k=3$}
Substituting the fixed index into the entry's wording gives an ordinary one-parameter array
count, a walk on $S=452$ vertices. Berlekamp--Massey on $2S=904$ exact terms
returns a minimal recurrence of order $38$ --- the order the entry states --- with
integer coefficients
\[
(c_1,\dots,c_{38})=(149,\ -7971,\ 165544,\ 514201,\ -80314173,\ 1119086138,\ 2353971653,\ -184820183374,\ 1195717243935,\ 7121767253646,\ -109236673637011,\ 151091595287238,\ 3146739450435147,\ -12340993650546304,\ -35536984802131697,\ 240979612131734495,\ 151218340006622456,\ -2409462261147296355,\ 214663115768358527,\ 14616659743291780241,\ -4465686520136547226,\ -57198163922758996002,\ 17039320101208544920,\ 146671703056688867968,\ -28570447814109231520,\ -245285612789187739552,\ 15003016836137771136,\ 261875539621553974784,\ 18764289131957794816,\ -170469142709804072960,\ -31680390599580352512,\ 61954278384892772352,\ 16293323636687241216,\ -10908963608341774336,\ -2882631151455830016,\ 965960077255114752,\ 165159806271750144,\ -40928839307624448).
\]
so that $a(m)=\sum_{i=1}^{38}c_i\,a(m-i)$ for every $m$. Removing the first
$38+4$ terms and repeating gives order $38$ again, which is what the
identification below needs.

\begin{theorem}
For each line above, the displayed recurrence holds for every index, and it is the recurrence
the entry records.
\end{theorem}

\begin{proof}
It holds by Lemma~\ref{lem:bm}. For the identification, the entry asserts that the line
satisfies SOME recurrence of the stated order, possibly only from some index on. Let $q$ be its
characteristic polynomial and $q_{\min}$ the minimal polynomial of the tail. Then
$q_{\min}\mid q$, and the second Berlekamp--Massey run --- on the line with its first few
terms removed --- shows $\deg q_{\min}$ equals the stated order, which is $\deg q$. Both are
monic, so $q=q_{\min}$.
\end{proof}

\section{Verification}

Three checks.

First, each line's model was compared against the entry's own published values for that line,
taken from the antidiagonal DATA read in either orientation and from the ``Table starts''
block, and agrees at every available term. This is the only point at which reading the entry's
English enters, and a misreading gives different counts.

Second, each recovered recurrence was evaluated directly on the model's values, with no
Berlekamp--Massey involved, and holds at every index tested.

Third, each order was computed twice by independent routes --- modulo a $61$-bit prime, where
every intermediate is a machine word, and again in exact rational arithmetic --- and once more
on the line with its first few terms dropped, which is what the identification requires. All
agree.

\begin{thebibliography}{9}
\bibitem{oeis} The OEIS Foundation, \emph{The On-Line Encyclopedia of Integer
Sequences}, \url{https://oeis.org}, 2026. Sequence A237158.
\bibitem{massey} J.~L.~Massey, Shift-register synthesis and BCH decoding, \emph{IEEE
Trans. Inform. Theory} \textbf{15} (1969), 122--127.
\bibitem{stanley} R.~P.~Stanley, \emph{Enumerative Combinatorics}, Volume 1, 2nd ed.,
Cambridge University Press, 2011. (Transfer-matrix method, Section 4.7.)
\bibitem{horn} R.~A.~Horn and C.~R.~Johnson, \emph{Matrix Analysis}, 2nd ed., Cambridge
University Press, 2013.
\end{thebibliography}

\end{document}
